\documentclass[11pt]{article}
\usepackage[margin=1in]{geometry}
\usepackage[T1]{fontenc}
\usepackage{lmodern}
\usepackage{microtype}

\begin{document}

\section{Limitations}
\label{sec:limitations}

The results should be interpreted in light of the implementation's scope and
the modeling assumptions documented in \texttt{agv\_core.py} and
\texttt{agv\_hybrid.py}. Some are deliberate simplifications that make the
planning procedure testable; others are limitations of the current prototype
that constrain its transfer to operational factories.

First, this is a simulation implementation. It has no robot operating system
(ROS) integration, vehicle localization or control interface, industrial
communications layer, or hardware-in-the-loop evaluation. This is a known
implementation limitation: simulation verification does not establish safe
execution under sensor error, vehicle dynamics, communication delay, or
controller behavior.

Second, each scenario supplies a static set of jobs known before planning.
The planner sorts jobs by desired pickup time and does not handle new arrivals
during execution. This is a deliberate batch-scheduling assumption, but it
limits conclusions about online manufacturing where urgent jobs, cancellations,
or changed deadlines require replanning. The fleet priority order is also
fixed (AGV-01 as $V_1$ through AGV-20 as $V_{20}$); priorities are not updated
to reflect waiting time, task urgency, or changing vehicle state. This
deliberate rule simplifies deterministic comparisons, while potentially
creating unfair delays and making performance dependent on the chosen order.

Third, every generated job is assigned and served exactly once. There is no
job reassignment, task preemption, vehicle breakdown model, or recovery
procedure after a failure. Exactly-once service is part of the modeled
problem, whereas the absence of failure handling is a known limitation for
deployment in factories where vehicles and stations may become unavailable.

Fourth, the parking node $\pi$ is deliberately exempt from node-conflict
checks so multiple AGVs may start and wait there. In the hybrid configuration,
lift lobbies can also be exempt waiting locations; lobby construction is
optional. The model does not impose a finite capacity on these exempt waiting
areas. Real parking bays and lift lobbies may have limited space, so the
assumption can understate congestion and should be revisited for each site.

Fifth, the implementation treats a physical track as exclusive in both
directions: it checks a traversal against reservations for either direction.
This is deliberately stricter than the paper's reverse-edge conflict check
and can reject safe same-direction following schedules. It favors
conservative separation but may overestimate congestion and completion time
relative to facilities that permit controlled following.

Sixth, experiments use one hand-crafted factory graph and one synthetic job
generator. This is a known generalizability limitation: topology, station
placement, task mix, and due-date distributions may differ substantially in
other plants. Finally, travel and lift times are deterministic functions of
edge lengths and speed, with no uncertainty, acceleration, or congestion
variation. The model also omits battery state and charging decisions. These
are deliberate scope choices for isolating route conflicts, but actual
travel-time variation and charging constraints can change feasible routes,
resource contention, and the relative performance of the architectures.

\end{document}
